% !TEX root = ../main.tex
\section{问题定义与实验设计}

\subsection{一般网络动力学问题}

设复杂系统包含 $N$ 个可观测组分，连接关系由邻接矩阵 $A$ 描述，全部组分在时刻 $t$ 的状态记为 $\mathbf X(t)=[\mathbf x_1(t),\ldots,\mathbf x_N(t)]$。
离散时间网络动力学可以一般地写为
\begin{equation}
\mathbf X(t+1)=\mathcal F_{\theta}\!\left(\mathbf X(t),A\right)+\boldsymbol\varepsilon(t),
\label{eq:network-dynamics}
\end{equation}
其中 $\mathcal F_{\theta}$ 表示未知或部分未知的演化规律，$\boldsymbol\varepsilon$ 概括过程随机性与观测误差。
若相互作用具有局部性和共享性，
节点更新还可写成
\begin{equation}
\mathbf x_i(t+1)\sim
f_{\theta}\!\left(
\mathbf x_i(t),+\operatorname{Agg}_{j:A_{ij}=1}
g_{\theta}(\mathbf x_i(t),\mathbf x_j(t))
\right).
\label{eq:local-interaction}
\end{equation}
这一表达将影响节点下一状态的因素分为两类：节点自身当前状态，以及由网络传入的邻居作用。
本文的方法围绕这一区分展开，而不要求 $f_{\theta}$ 具有某个特定动力学的解析形式。

据此定义两个递进任务。
网络已知时，给定 $A$ 和连续观测到的状态转移对，学习 $f_{\theta}$，并检验模型能否在未显式写入方程的情况下恢复正确的局部响应。
网络未知时，
只给出状态序列，
同时估计动力学参数 $\theta$ 与潜在邻接矩阵 $\hat A$。
第二项任务不仅要求预测准确，
还要求输出的连接能够在边级评价中接近真实网络。

\subsection{网络 SIR 动力学}

本文用离散时间网络 SIR 过程具体化\cref{eq:network-dynamics}。
在无向无自环网络上，节点 $i$ 的状态 $s_i(t)\in\{S,I,R\}$ 分别表示易感、感染和康复。
感染节点以概率 $\gamma$ 康复；
易感节点受到其感染邻居的独立影响。
时刻 $t$ 的感染者邻居数为
\begin{equation}
k_i(t)=\sum_j A_{ij}\,\mathbf{1}[s_j(t)=I],
\end{equation}
其中，$\mathbf{1}[\cdot]$ 为示性函数：方括号内的条件成立时取 1，否则取 0。
因此，$\mathbf{1}[s_j(t)=I]$ 仅在节点 $j$ 处于感染态时取 1，
与邻接矩阵元素 $A_{ij}$ 相乘并求和后，即得到节点 $i$ 的感染者邻居数。
相应的一步感染概率为
\begin{equation}
P(S\to I\mid k_i)=1-(1-\beta)^{k_i},
\label{eq:sir-infection}
\end{equation}
其中 $\beta$ 是单条感染边上的传播概率。
SIR 对本文尤其合适：一方面，节点自身状态决定可发生的转移，邻居计数决定交互强度，
与双通道思想直接对应；
另一方面，
\cref{eq:sir-infection}给出了可检验的真实响应曲线，
使“预测看起来正确”与“动力学规律确实被恢复”能够被区分。

\subsection{实验设计}

实验使用三类具有代表性的合成网络：ER 随机网络、WS 小世界网络和 BA 无标度网络。
每类只生成一个固定的 $N=100$ 图实例，
统一设置 $\beta=0.3$、$\gamma=0.1$。
对多组初始感染节点和随机传播轨迹进行模拟，
最终得到 9000 个相邻时刻状态转移对。
主要设置见\cref{tab:protocol}。

\begin{table}[H]
  \centering
  \caption{SIR 验证场景的数据与评价设置}
  \label{tab:protocol}
  \begin{tabularx}{0.96\textwidth}{@{}lX@{}}
    \toprule
    项目 & 设置 \\
    \midrule
    网络 & ER$(N=100,p=0.1)$、WS$(N=100,k=6,p=0.2)$、BA$(N=100,m=3)$，各一个固定实例 \\
    动力学 & 离散同步 SIR，$\beta=0.3$，$\gamma=0.1$，每条轨迹 60 个状态转移 \\
    观测数据 & 30 次初始节点抽样，每次模拟 5 条随机轨迹，共 9000 个状态转移对 \\
    数据划分 & 状态转移对按 $8:2$ 随机划分为训练集和验证集 \\
    规律诊断 & 感染概率随邻居数 $k$ 的变化、$\beta$ 与 $\gamma$ 的参数反演 \\
    多步推演 & 5 组初始条件、1--5 个初始感染者，比较 S/I/R 占比轨迹 \\
    结构评价 & 硬图的 ACC、Precision、F1、TPR、FPR 及预测边数 \\
    \bottomrule
  \end{tabularx}
\end{table}

评价遵循由“预测”到“规律”再到“结构”的层次。
一步预测使用逐节点交叉熵与准确率；
规律诊断比较模型给出的 $P(S\to I\mid k)$ 与理论曲线，
并从输出曲线反演 $\hat\beta$，
同时考察 $P(I\to R)$ 得到的 $\hat\gamma$；
多步推演把模型输出重新作为下一步输入，
统计群体 S/I/R 占比的平均绝对误差。
未知网络还使用
\begin{equation}
\mathrm{Precision}=\frac{TP}{TP+FP},\qquad
F_1=\frac{2\,\mathrm{Precision}\,\mathrm{TPR}}
{\mathrm{Precision}+\mathrm{TPR}}
\end{equation}
评价恢复的具体边，避免仅由大量真负边造成的高 ACC 掩盖稠密错误图。

多次训练采用同一图和同一数据上的不同优化初始化。
具体训练配置、结果文件与指标计算方式统一列于附录。
